\subsection*{一步一步建立直观}

\paragraph{1. 把乘以虚数看成速度规则}

对复平面上的位置 $z$，乘以 $i$ 会把它逆时针转过直角。如果规定速度为 $z'=iz$，速度就始终与半径垂直，且大小与半径相等。右图的蓝色短箭头表示这个方向，位置点因此沿圆运动，而不是沿固定直线前进。



\paragraph{2. 连续执行局部规则得到整体旋转}

转到半圈时，单位初值到达负实轴。复指数把整段演化记成 $e^{it}$：指数中的参数是持续执行局部旋转规则的时间。它与实指数的增长共享同一种微分结构，只是实数倍率改变长度，虚数倍率改变方向。



\paragraph{3. 实部与虚部各负责一种变化}

加入正的实部增长率后，速度既有径向分量，也有切向分量，轨迹变为螺线。公式 $e^{(\sigma+i)t}=e^{\sigma t}e^{it}$ 正好把这两种作用分开。窗口固定角速度为一；一般角速度应写入虚部，并与实际时间单位保持一致。


